\documentclass[11pt]{article}
\usepackage{amsmath,amssymb,amsthm,geometry,hyperref}
\hypersetup{hidelinks,pdftitle={A sharp spectral-gap refinement for the switch chain},pdfauthor={Tom Ayre}}
\setlength{\emergencystretch}{2em}
\geometry{margin=1in}
\newtheorem{theorem}{Theorem}
\newcommand{\Var}{\operatorname{Var}}
\newcommand{\gap}{\operatorname{gap}}
\title{A sharp spectral-gap refinement for the switch chain}
\author{Tom Ayre\\\small ayre.tom at gmail.com}
\date{October 7, 2026}
\begin{document}\maketitle
\paragraph{Status.} Conditional research note prepared with assistance from Codex. This note retains the
triple-resampling inequality and slice equality-projection inequality of
OpenAI's \emph{Polynomial Mixing of the Switch Chain for Every Graphical
Degree Sequence} (September 25, 2026). The new step strengthens its
nonnegative disjoint-pair estimate to $D\succeq H$. We also use the exact
Johnson-slice gap. A targeted literature search did not locate this
statement; publication priority is not established. The source was read at
the \href{https://github.com/openai/math/tree/adc7f1241b42e322a6451854ab7e4b4c146bf78a/preprints/Polynomial-Mixing-of-the-Switch-Chain-for-Every-Graphical-Degree-Sequence-September-25-2026}{archived source directory} at revision\\
\texttt{adc7f1241b42e322a6451854ab7e4b4c146bf78a}.
Separate agent reviews of the source's finite row, triple, slice and
disjoint-cell arguments found no gap. These are local AI-assisted checks,
not formal verification or independent external human peer review.

Let $\Omega_d$ be the simple labeled graphs with fixed graphical degree
sequence $d$ on $n\ge4$ vertices. For a vertex pair $a$, let $E_a$ be
conditional uniform resampling of its singleton neighbors, $h_a=I-E_a$,
and $H=\sum_a h_a$. Let $L$ be the unnormalized Laplacian whose states
are $\Omega_d$ and whose edges are single valid switches, each neighbor
counted once. Set $D=\sum_{a\cap b=\varnothing}h_a h_b$, summing ordered
pairs. All inner products use the uniform probability measure.

\begin{theorem}[Conditional candidate]
Assume the source's triple inequality $H_T^2\succeq H_T$ and slice
inequality $S\preceq G$ described below. Then
\[
 D\succeq H,\qquad H^2\succeq2H,\qquad H\preceq L.
\]
If $|\Omega_d|>1$, then $\gap(L)\ge2$, with sharp universal constant.
For the source's lazy four-vertex-proposal chain,
\[
 \gap(P_d)\ge\frac1{6\binom n4},\qquad
 t_{\rm mix}(\varepsilon)\le
 \left\lceil6\binom n4\left(\tfrac12\log|\Omega_d|
                  +\log\frac1{2\varepsilon}\right)\right\rceil
 \quad(0<\varepsilon<1/2).
\]
Thus $t_{\rm mix}(1/4)=O(n^6)$ rather than the displayed $O(n^8)$ bound.
The spectral-gap bound is attained for $d=(2,2,1,1,0,\ldots,0)$.
The nonlazy uniform-pair heat-bath chain $K=\binom n2^{-1}\sum_a E_a$
has sharp universal gap $\gap(K)\ge2/\binom n2$.
\end{theorem}

\paragraph{Proposal conventions.}
The chain $P_d$ holds with probability $1/2$. Otherwise it chooses a
four-vertex set uniformly, chooses uniformly one of the six ordered
pairs of distinct perfect matchings on that set, and replaces the first
matching by the second if all removed edges are present and all added
edges are absent; otherwise it holds. Thus every distinct switch
neighbor has transition probability $1/[12\binom n4]$.
The chain $K$ chooses a pair uniformly and samples uniformly from its
entire pair fiber, including its current state. This chain is a convex
combination of orthogonal projections, so its eigenvalues are nonnegative.

\paragraph{Borrowed inputs.}
For disjoint pairs $a,b$, let $g_{a,b}=(I-T_{a,b})/2$, where $T_{a,b}$
switches the four cross edges when possible. On an $a$-fiber, $R_{a,b}$
projects onto the centered indicator that the two singleton neighbors in
$b$ receive the same endpoint. It is zero if either is not a singleton,
or if the indicator is constant. On a uniform $q$-subset slice of $N$
singleton labels put
\[
 G=\tfrac12\sum_{k<l}(I-\tau_{kl}),\qquad S=\sum_{k<l}Q_{kl},
\]
where $Q_{kl}$ is the equality-indicator projection. Source Lemma
\texttt{sl:comparison} proves $S\preceq G$. The elementary Johnson calculation in the appendix
gives $G\succeq (N/2)(I-E)$ on nonconstant slices. Source Proposition
\texttt{tri:gap} proves $H_T^2\succeq H_T$. The disjoint-cell decomposition
used next is spelled out in its \texttt{sections/disjoint.tex}.

\paragraph{A sharper disjoint-cell estimate.}
Fix disjoint pairs $a,b$. In a source disjoint cell, outside singleton
sets have sizes $A,B$; the $2\times2$ cross matrix has total $s$, first
row sum $r$ and first column sum $c$. The counts of compatible outside
subsets are $a_r=\binom A{k+1-r}$ and $b_c=\binom B{l+1-c}$.
In the dependent case $s=2$, the matrix-multiplicity table has center
two, four adjacent entries one, and zero corners. If both off-center
events have positive mass, $a_1b_1>0$. Set
\[
 \alpha=\frac{a_0+a_2}{a_1},\quad
 \beta=\frac{b_0+b_2}{b_1},\quad
 \rho^2=\frac{\alpha\beta}{(2+\alpha)(2+\beta)}.
\]
Here $\rho$ is the magnitude of the rank-one correlation in the source.
For $0\le k\le A$,
\[
 \alpha=\frac{A-k}{k+1}+\frac{k}{A-k+1}\le(A-k)+k=A,
\]
and similarly $\beta\le B$. Thus, for $\kappa_A=A/(A+2)$ and
$\kappa_B=B/(B+2)$, we have $\rho^2\le\kappa_A\kappa_B$.

The source's correlation-plane component $f=xu+yv$ satisfies
$h_a f=y(v+\rho u)$, $h_b f=x(u+\rho v)$; each parenthesized vector has
squared norm $1-\rho^2$ and their mutual product is $\rho(1-\rho^2)$.
Weighted Young gives
\[
 2xy\rho(1-\rho^2)\ge
 -\kappa_A y^2(1-\rho^2)-\kappa_B x^2(1-\rho^2).
\]
The first error line lies in $\operatorname{ran}R_{a,b}$ and the second
in $\operatorname{ran}R_{b,a}$, as in the source. The common orthogonal
component contributes at least $2\|g_{a,b}f\|^2$. Therefore
\[
 2\langle h_af,h_bf\rangle\ge2\|g_{a,b}f\|^2
       -\kappa_A\|R_{a,b}f\|^2-\kappa_B\|R_{b,a}f\|^2.\tag{*}
\]
Independent cells satisfy this without the error terms. The coefficients
have a fiberwise interpretation: $R_{a,b}=0$ unless $c=1$, and on such
fibers the singleton count is $N=A+2$. Its coefficient is therefore
$1-2/N$; symmetrically for $R_{b,a}$. Constant slices have $h_a=0$.

\paragraph{Summing the positive remainder.}
Sum (*) over unordered disjoint pairs. For fixed $a$, the resulting
operator on a nonconstant fiber is
\[
 G-(1-2/N)S\succeq(2/N)G\succeq I-E.
\]
The first step uses $S\preceq G$ and $N\ge2$, the second the Johnson
gap. Averaging and summing over $a$ gives $D\succeq H$.
\newpage
The source's exact identity is
\[
 H^2=\sum_{T\in\binom{[n]}3}H_T^2-(n-3)H+D.
\]
The triple inequality now gives $H^2\succeq2H$. Connectivity by switches
identifies its kernel with constants, hence
$2\Var(f)\le\langle f,Hf\rangle$.
Let $J_a=\sum_{k<l}(I-\tau_{kl})$ on each fiber. Its nonzero gap is $N$,
so $h_a\preceq J_a/N\preceq J_a/2$. Each switch has exactly two pair
representations; thus $\sum_aJ_a=2L$, and $H\preceq L$.
This proves $\gap(L)\ge2$. The source normalization is
$I-P_d=L/[12\binom n4]$; the usual $L^2$-to-TV estimate gives the
stated mixing bound.

\paragraph{Sharpness.}
For $d=(2,2,1,1)$ the only realizations have edge sets $\{12,13,24\}$
and $\{12,14,23\}$. They are joined by one switch; the Laplacian has
spectrum $\{0,2\}$. Adding isolated vertices preserves this state space.
The lazy proposal kernel has off-diagonal entry $1/[12\binom n4]$,
so its gap is exactly $1/[6\binom n4]$.
Exactly two vertex pairs have nonconstant fibers in this example.
Each such resampling is averaging on the two states, so $H$ has spectrum
$\{0,2\}$ and $\gap(K)=2/\binom n2$ as well.

\paragraph{Checks and limitations.}
The operator diagnostic checks $D\succeq H$, $H^2\succeq2H$, $H\preceq L$
and the nontrivial spectral gaps for all 144 degree-sequence classes on
$n=4,5,6$, scanning 33,856 labeled graphs. A separate implementation
constructs $L$ directly from switches, without using pair fibers or the
operators $H,D$. It checks all 486 degree-sequence classes on $n=4,5,6,7$,
including 366 nontrivial classes. The minimum observed gap is 2, within
tolerance $10^{-7}$. These are finite floating-point diagnostics, not proofs.
The source's triple and equality-projection inequalities were reviewed
locally in separate Codex passes without locating a gap. External peer
review and formal verification remain outstanding; no fresh Lean kernel
verification was performed. No definitive novelty claim is made. Only
the spectral-gap constant is asserted sharp, not the $O(n^6)$ mixing exponent.

\paragraph{Primary literature checked.}
The closest located sharp analogue is Fu, Qin and Wang,
\emph{Spectral Gap for the Binary Fixed-Margin Swap Chain},
\url{https://arxiv.org/abs/2606.22636}. It treats binary fixed-margin
matrices, equivalently bipartite graphs, rather than the general undirected
state space considered here. It also supplies the three-row inequality
used in the source framework.
Johnson-spectrum switch/Curveball comparison is already standard:
\url{https://arxiv.org/abs/1709.07290} and
\url{https://drops.dagstuhl.de/entities/document/10.4230/LIPIcs.APPROX-RANDOM.2018.36}.
For undirected Curveball and the two pair representations per switch,
see \url{https://pmc.ncbi.nlm.nih.gov/articles/PMC6072652/}.
For the unified switch-chain analysis see
\url{https://arxiv.org/abs/1903.06600}.
The potentially new part is the sharper disjoint-interaction estimate
and sharp universal spectral-gap consequence within the source framework.
A targeted search found no matching general undirected statement; it does
not rule out a known implicit consequence or establish publication priority.
\clearpage
\appendix
\section{An elementary Johnson-slice gap proof}
Let $\mathcal X_{N,q}$ be the uniform $q$-subsets of $[N]$, with
$1\le q\le N-1$, and let $J=\sum_{i<j}(I-\tau_{ij})$.
For $A\subseteq[N]$ put $x_A=\prod_{i\in A}x_i$. Let $V_k$ be the span
of the monomials of degree at most $k$ on the slice. For $|A|=k$, direct
summation over transpositions with one endpoint in $A$ gives
\[
 Jx_A=k(N-k+1)x_A-(q-k+1)\sum_{a\in A}x_{A\setminus\{a\}}.\tag{J}
\]
Indeed before using $\sum_i x_i=q$ the expression is
$k(N-k)x_A-\sum_{a\in A}x_{A\setminus\{a\}}\sum_{b\notin A}x_b$,
and each product in the second sum equals
$(q-k+1)x_{A\setminus\{a\}}-x_A$.
Thus $J$ preserves $V_k$, and modulo $V_{k-1}$ acts as
$\lambda_k=k(N-k+1)$. Self-adjointness shows that on
$V_k\cap V_{k-1}^{\perp}$ it acts by exactly this scalar: the difference
$(J-\lambda_k I)f$ lies in both $V_{k-1}$ and its orthogonal complement.
Finally $V_q$ is the full function space, because degree-$q$ monomials
are point indicators. On each nonconstant nonzero level,
\[
 \lambda_k-N=(k-1)(N-k)\ge0.
\]
This proves $J\succeq N(I-E)$, or $G=J/2\succeq(N/2)(I-E)$.
The levels $q=0,N$ are constant and need no gap estimate.
\end{document}
